\Needspace{12\baselineskip}
\section{Rounding a fractional cover}\label{sec:rounding}

We convert the color bound into independent numerical variables on the
vertices. Their means will be small, while the common member supplied by
Theorem~\ref{thm:multiscale-selector} will force a positive amount of
fractional-cover mass onto sets whose products are at least one. Comparing
the resulting lower and upper bounds on the same expectation completes
the proof. This conversion follows the selector-to-rounding approach
of Dubroff, Kahn, and Park~\cite[Section~2]{DKP2024} and
Pham~\cite[arXiv version~1, Section~3]{Pham2024}. In particular, the
equal distribution of a set's weight among its vertices in
\eqref{round:mass} is Pham's construction. The simultaneous color
bound supplies the additional control used here for arbitrary set sizes.

\begin{proof}[Proof of Theorem~\ref{thm:main}]
Put
\[
 D=256,\qquad B=512,\qquad C=25B^4.
\]
Let $r\in(0,1]$ and suppose that $g:2^X\to[0,1]$ satisfies
\begin{equation}\label{round:cover}
 \sum_{S\subseteq H}g(S)\ge1\quad(H\in\mathcal F),
 \qquad
 \sum_{S\subseteq X}g(S)r^{|S|}\le\frac12.
\end{equation}
Set $p=r/C$. We prove that $\mathcal F$ is $p$-small.

First observe that $g(\varnothing)\le1/2$, since its contribution to the
second sum in \eqref{round:cover} is $g(\varnothing)$. For each
$H\in\mathcal F$, define
\begin{equation}\label{round:mass}
 M_H=\sum_{\varnothing\ne S\subseteq H}g(S)
 \ge1-g(\varnothing)\ge\frac12,
 \qquad
 \lambda_H(x)=\frac1{M_H}
 \sum_{\substack{\varnothing\ne S\subseteq H\\x\in S}}
 \frac{g(S)}{|S|}\quad(x\in X).
\end{equation}
The function $\lambda_H$ vanishes off $H$ and is a probability mass
function: summing over $x$ counts each nonempty $S$ exactly $|S|$ times.

Suppose, for a contradiction, that $\mathcal F$ is not $p$-small.
Let $s$ be the largest nonnegative integer for which
\begin{equation}\label{round:scale}
 p\sum_{i=1}^sD^i\le\frac12.
\end{equation}
Such an integer exists: zero satisfies the restriction, and $p>0$ makes
the geometric sums unbounded. Apply
Theorem~\ref{thm:multiscale-selector} to these probability masses and this
choice of $s$. Thus the colors $a(x)$ are independent, with
\[
 \P(a(x)=i)=\pi_i=D^ip\quad(1\le i\le s),
 \qquad
 \P(a(x)=s+1)=\pi_{s+1}=1-p\sum_{i=1}^sD^i.
\]
With probability at least $9/10$, there is an $H\in\mathcal F$ such that
$\sum_x a(x)\lambda_H(x)\le2$. Denote this event by $\mathcal E$, and set
\begin{equation}\label{round:variables}
 Y_x=B^{4-a(x)},\qquad
 Z=\sum_{\varnothing\ne S\subseteq X}g(S)\prod_{x\in S}Y_x.
\end{equation}

Fix an outcome in $\mathcal E$ and a corresponding member $H$. Writing
$\overline a(S)=|S|^{-1}\sum_{x\in S}a(x)$ for nonempty $S$, we have
\[
 \frac1{M_H}\sum_{\varnothing\ne S\subseteq H}
 g(S)\overline a(S)
 =\sum_x a(x)\lambda_H(x)\le2.
\]
Apply Markov's inequality to the probability distribution
$g(S)/M_H$ on the nonempty subsets of $H$. It gives
\[
 \sum_{\substack{\varnothing\ne S\subseteq H\\
                  \overline a(S)\le4}}g(S)\ge\frac{M_H}{2}.
\]
For every set in this sum,
\[
 \prod_{x\in S}Y_x
 =B^{\,|S|(4-\overline a(S))}\ge1.
\]
Consequently $Z\ge M_H/2\ge1/4$ on $\mathcal E$. This holds for each
outcome, although the corresponding $H$ may depend on the coloring.
Since $Z$ is nonnegative everywhere,
\begin{equation}\label{round:lower}
 \E Z\ge\frac14\P(\mathcal E)\ge\frac9{40}.
\end{equation}

For the upper bound, maximality of $s$ gives
\[
 \frac12<p\sum_{i=1}^{s+1}D^i\le2pD^{s+1}.
\]
As $B\ge D$, it follows that
\begin{equation}\label{round:tail}
 B^{-(s+1)}\le D^{-(s+1)}<4p.
\end{equation}
The variables $Y_x$ are independent and identically distributed, and
their common mean $\mu$ satisfies
\begin{align}
 \mu
 &=B^4\left(p\sum_{i=1}^s(D/B)^i
            +\pi_{s+1}B^{-(s+1)}\right)\notag\\
 &\le B^4\left(p\sum_{i=1}^s2^{-i}+4p\right)
 \le5B^4p=\frac r5.\label{round:mean}
\end{align}
The same estimates hold when $s=0$, with the first sums empty.
Independence now yields
\begin{align}
 \E Z
 &=\sum_{\varnothing\ne S\subseteq X}g(S)\mu^{|S|}\notag\\
 &\le\sum_{\varnothing\ne S\subseteq X}g(S)(r/5)^{|S|}
 \le\frac15\sum_{\varnothing\ne S\subseteq X}g(S)r^{|S|}
 \le\frac1{10}.\label{round:upper}
\end{align}
Here the factor $1/5$ is valid because every set in the sum is
nonempty. The bounds \eqref{round:lower} and \eqref{round:upper}
contradict one another.

Thus $\mathcal F$ is $r/C$-small for every admissible $r>0$. At $r=0$,
the family $\mathcal F$ itself is a cover of cost zero, since
$\varnothing\notin\mathcal F$. Taking $r=q_f(\mathcal F)$, with this
last observation covering the case $q_f(\mathcal F)=0$, proves
\[
 q_f(\mathcal F)\le25\cdot512^4\,q(\mathcal F).
\]
\end{proof}
